% Section 10 ('Appendix 10' in the paper's cross-references): D12 -- assignment of the in-plane harmonics e^{+-i m phi} to irreps by m mod 12 (Eq. 24), enforced nodes, Sym^2 E_m
% residue 0 stands for m = 12, 24, ...: the constant m = 0 alone is A_1, A_2 (sin m phi) first appears at m = 12
\begin{tabular}{ll}
\hline
$m \bmod 12$ & irreps \\
\hline
0 & $A_1 \oplus A_2\ (m \ge 12;\ m = 0:\ A_1)$ \\
1 & $E_1$ \\
2 & $E_2$ \\
3 & $E_3$ \\
4 & $E_4$ \\
5 & $E_5$ \\
6 & $B_1 \oplus B_2$ \\
7 & $E_5$ \\
8 & $E_4$ \\
9 & $E_3$ \\
10 & $E_2$ \\
11 & $E_1$ \\
\hline
\end{tabular}

\begin{tabular}{lccc}
\hline
irrep & node on $C_2'$ & node on $C_2''$ & node on the 12-fold axis \\
\hline
$A_1$ & - & - & - \\
$A_2$ & node & node & - \\
$B_1$ & - & node & node \\
$B_2$ & node & - & node \\
$E_1$ & - & - & node \\
$E_2$ & - & - & node \\
$E_3$ & - & - & node \\
$E_4$ & - & - & node \\
$E_5$ & - & - & node \\
\hline
\end{tabular}

\begin{tabular}{llc}
\hline
$E_m$ & $\mathrm{Sym}^2 E_m$ & quartic invariants \\
\hline
$E_1$ & $A_1 \oplus E_2$ & 2 \\
$E_2$ & $A_1 \oplus E_4$ & 2 \\
$E_3$ & $A_1 \oplus B_1 \oplus B_2$ & 3 \\
$E_4$ & $A_1 \oplus E_4$ & 2 \\
$E_5$ & $A_1 \oplus E_2$ & 2 \\
\hline
\end{tabular}

% weak coupling on the circle: R = 3/2 (real), 1 (chiral)
